\chapter{Minacce alla validità}
\label{cap:validita}

Le minacce sono classificate secondo Wohlin et al.: per ognuna si indica la mitigazione adottata. Quelle non eliminabili sono dichiarate e pesano sull'interpretazione dei risultati.

\section{Validità interna}
\textit{La differenza misurata è causata dal fattore dichiarato?}

\begin{tabularx}{\linewidth}{@{}LL@{}}
\toprule
\textbf{Minaccia} & \textbf{Mitigazione} \\
\midrule
Interferenza dell'host: generatore di carico, cluster e sistema operativo condividono CPU e memoria & Soglia di CPU, nessuno swap ammesso, carichi sotto la capacità misurata, run non validi scartati \\
Stato trascinato tra run (backoff dei riavvii, tabella che cresce, worker scalati, topologia di V4 dopo un failover) & Reset che ricrea i pod e verifica lo stato iniziale prima di partire \\
Degrado progressivo dell'ambiente durante la campagna & Ordine dei run casuale \\
Istante del guasto impreciso & Registrato nello stesso comando che invia il segnale, su entrambi gli orologi \\
Proxy di iniezione che altera la latenza solo in alcuni run & Toxiproxy sempre nel percorso verso il database in V2 e V3 \\
\bottomrule
\end{tabularx}

\section{Validità di costrutto}
\textit{Misuriamo davvero ciò che intendiamo misurare?}

\begin{tabularx}{\linewidth}{@{}LL@{}}
\toprule
\textbf{Minaccia} & \textbf{Mitigazione} \\
\midrule
La soglia del 99\% per secondo dipende dal ritmo: a 50 richieste/s basta un errore per perdere il secondo & Si riportano sempre anche le richieste fallite per guasto \\
La finestra di 5 s di ``ripristinato'' è arbitraria & Analisi di sensibilità con 3 s e 10 s \\
Il tempo end-to-end esclude l'attesa del commit & Attesa visibile nella latenza lato client; scelta dichiarata in D0005 \\
\texttt{kubectl delete pod} non rappresenta un guasto & Crash e blocco con segnali dal nodo (D0016) \\
\bottomrule
\end{tabularx}

\section{Validità esterna}
\textit{Quanto si generalizzano i risultati?}

\begin{tabularx}{\linewidth}{@{}LL@{}}
\toprule
\textbf{Minaccia} & \textbf{Mitigazione} \\
\midrule
In locale le repliche girano sulla stessa macchina fisica: ridondanza di processo, non di hardware & Dichiarato; la fase AWS usa nodi separati (RQ6) \\
Il tempo di servizio è simulato come attesa, non come calcolo: più worker producono più capacità anche su un solo host & Dichiarato in D0015; i risultati di V5 valgono per carichi limitati dall'I/O \\
Servizio giocattolo, un solo tipo di richiesta & Le conclusioni riguardano i meccanismi (rilevamento, failover, smaltimento), non i valori assoluti \\
kind non è un cluster di produzione & Confronto con EKS (RQ6) \\
\bottomrule
\end{tabularx}

\section{Validità delle conclusioni}
\textit{Le conclusioni reggono statisticamente?}

\begin{tabularx}{\linewidth}{@{}LL@{}}
\toprule
\textbf{Minaccia} & \textbf{Mitigazione} \\
\midrule
Poche ripetizioni & 10 ripetizioni per le configurazioni decisive, intervalli bootstrap, differenza minima rilevabile dichiarata \\
Interpretazione adattata ai dati & Previsioni committate prima della campagna \\
Molti confronti, qualche ``differenza'' per caso & Si confrontano solo le coppie previste dalle domande di ricerca \\
\bottomrule
\end{tabularx}
